\chapter{小数的运算}\label{ch:g5:decimalops}

小数的加法、减法和乘法与整数几乎一样，
关键在于知道小数点该放在哪里。本章
介绍竖式计算方法，以及乘以或除以
$10$、$100$、$1000$ 的妙处。（两个小数相乘
见\cref{ch:g6:decimals}。）

\section{加法和减法}

\begin{method}[小数点对齐的竖式]\label{met:g5:decimalops:addsub}
\begin{enumerate}
  \item 把各数写成\emph{小数点对齐}的形式
        （个位对个位，十分位对十分位）；
  \item 在位数较少的小数部分末尾补零；
  \item 像整数一样进行加法或减法；
  \item 把小数点直接向下写到结果里。
\end{enumerate}
\end{method}

\begin{example}\label{ex:g5:decimalops:addsub}
计算 $27.5 + 3.86$（补写 $27.5 = 27.50$）：
\[
\begin{array}{r}
2\,7.5\,0 \\
+\ \ \ 3.8\,6 \\
\hline
3\,1.3\,6
\end{array}
\qquad\qquad
\begin{array}{r}
1\,2.4\,0 \\
-\ \ \ 5.6\,5 \\
\hline
\ \ \,6.7\,5
\end{array}
\]
减法 $12.4 - 5.65$：补零，照常借位，
再用加法验算：$6.75 + 5.65 = 12.4$。
\end{example}

\begin{example}[求补数]\label{ex:g5:decimalops:complement}
要使 $7.6$ 变成 $10$，需要加上多少？先凑到整数：$7.6 \to 8$ 需要
$0.4$，再从 $8 \to 10$ 需要 $2$，所以补数是 $2.4$。找零时很有用：
用 $10$ 元买一件售价 $7.60$ 元的商品，应找回 $2.40$ 元。
\end{example}

\section{扩大十倍，缩小十倍}

\begin{proposition}[乘以和除以 10、100、1000]\label{prop:g5:decimalops:shift}
乘以 $10$ 时，每个数字所代表的值都扩大为原来的十倍：
各数字都向左移动一位，这看起来像是\emph{小数点
向右移动了一位}。除法则正好相反：
\[
3.47 \times 10 = 34.7, \qquad
3.47 \times 100 = 347, \qquad
52.1 \div 10 = 5.21, \qquad
6 \div 100 = 0.06 .
\]
\end{proposition}
\admitted

\begin{omfigure}
\begin{tikzpicture}[scale=0.9]
  \draw (0,0) grid[xstep=1.5, ystep=0.75] (6.0,2.25);
  \node[font=\small] at (0.75,2.6) {十位};
  \node[font=\small] at (2.25,2.6) {个位};
  \node[font=\small] at (3.75,2.6) {十分位};
  \node[font=\small] at (5.25,2.6) {百分位};
  \node[font=\normalsize, omExo] at (2.25,1.85) {$3$};
  \node[font=\normalsize, omExo] at (3.75,1.85) {$4$};
  \node[font=\normalsize, omExo] at (5.25,1.85) {$7$};
  \node[font=\normalsize, omDef] at (0.75,1.1) {$3$};
  \node[font=\normalsize, omDef] at (2.25,1.1) {$4$};
  \node[font=\normalsize, omDef] at (3.75,1.1) {$7$};
  \node[right, font=\small] at (6.2,1.85) {$3.47$};
  \node[right, font=\small, omDef] at (6.2,1.1)
    {$3.47 \times 10 = 34.7$};
  \draw[-Stealth, omDef, thick] (4.4,1.7) -- (3.1,1.25);
\end{tikzpicture}

{\small 乘以 $10$：每个数字都向左移动一位。
小数点其实并没有移动，移动的是数字。}
\end{omfigure}

\begin{example}[再谈单位换算]\label{ex:g5:decimalops:convert}
相邻单位之间的进率都是 10 的幂，因此
单位换算只需进行这样的移位：$2.35$ 米 $= 235$ 厘米（$\times 100$）；$470$ 克
$= 0.47$ 千克（$\div 1000$）；$1.5$ 升 $= 150$ 厘升。
\end{example}

\section{小数乘整数}

\begin{method}[小数乘整数]\label{met:g5:decimalops:mult}
先不看小数点，按整数相乘；再使积的小数位数
与该小数因数的小数位数一样多：
\[
4.35 \times 6: \quad 435 \times 6 = 2\,610
\ \to\ 4.35 \times 6 = 26.10 = 26.1 .
\]
先估算，以便有把握地确定小数点的位置：$4.35$ 接近
$4$，而 $4 \times 6 = 24$，所以结果是 $26.1$，而不是 $2.61$，也不是 $261$。
\end{method}

\begin{example}[重复加法依然适用]\label{ex:g5:decimalops:mult}
$0.75 \times 4 = 0.75 + 0.75 + 0.75 + 0.75 = 3$：四个四分之一
合成一个整体，只是披上了小数的外衣。乘以整数仍然意味着“
相同的数相加了多少次”。
\end{example}

\begin{example}[购物]\label{ex:g5:decimalops:shopping}
三本笔记本，每本 $2.45$，一支钢笔 $1.20$：
\begin{enumerate}
  \item 笔记本：$2.45 \times 3 = 7.35$；
  \item 总价：$7.35 + 1.20 = 8.55$；
  \item 付 $10$ 时的找零：$10 - 8.55 = 1.45$。
\end{enumerate}
\end{example}

\section{练习}

\begin{exercise}[$\star$]\label{exo:g5:decimalops:1}
用竖式计算：$45.7 + 8.93$；\quad $16.25 + 3.75$；\quad
$0.86 + 12.4$。
\end{exercise}

\begin{exercise}[$\star$]\label{exo:g5:decimalops:2}
用竖式计算并验算：$9.4 - 2.72$；\quad $20 - 13.45$；
\quad $6.03 - 5.9$。
\end{exercise}

\begin{exercise}[$\star$]\label{exo:g5:decimalops:3}
口算下列各数凑成 $10$ 所需的补数（见
\cref{ex:g5:decimalops:complement}）：$6.2$；\quad $3.75$；\quad
$9.99$。
\end{exercise}

\begin{exercise}[$\star$]\label{exo:g5:decimalops:4}
不用竖式，直接计算：
\[
7.24 \times 10, \qquad
0.58 \times 100, \qquad
34.5 \div 10, \qquad
7 \div 100, \qquad
0.3 \times 1000 .
\]
\end{exercise}

\begin{exercise}[$\star$]\label{exo:g5:decimalops:5}
换算下列单位：$4.05$ m 化为 cm；\quad $325$ cm 化为 m；\quad $0.6$ kg 化为 g；
\quad $85$ cL 化为 L。
\end{exercise}

\begin{exercise}[$\star$]\label{exo:g5:decimalops:6}
先估算，再计算：$6.2 \times 4$；\quad
$3.45 \times 8$；\quad $12.5 \times 6$。
\end{exercise}

\begin{exercise}[$\star$]\label{exo:g5:decimalops:7}
跑道一圈长 $0.4$ km。跑 $7$ 圈和
$25$ 圈各长多少？
\end{exercise}

\begin{exercise}[$\star$]\label{exo:g5:decimalops:8}
莉亚买了 $2$ 根法棍，每根 $1.15$，又买了一个蛋糕，价格 $12.60$。她
用一张 $20$ 的纸币付款。计算她的总价和找零。
\end{exercise}

\begin{exercise}[$\star$]\label{exo:g5:decimalops:9}
一瓶装 $1.5$ L。一包 $6$ 瓶共有多少升？
一家人每天喝 $0.75$ L，一包够喝
多少天？（$0.75$ 能在 $9$ 里除几次？）
\end{exercise}

\begin{exercise}[$\star\star$]\label{exo:g5:decimalops:10}
米洛算出 $5.6 \times 3 = 15.18$，他的推理是“$5 \times 3 = 15$
和 $6 \times 3 = 18$”。用估算说明答案一定错了，
再正确计算。
\end{exercise}

\begin{exercise}[$\star\star$]\label{exo:g5:decimalops:11}
一位跑步者跑四圈，用时分别为 $65.4$ s、$64.9$ s、$65.0$ s 和
$66.1$ s。
\begin{enumerate}
  \item 计算四圈的总用时。
  \item 四圈的纪录是 $260$ s。她是否打破了纪录？
        若打破，打破了多少秒？
\end{enumerate}
\end{exercise}
